


\section{Thiết kế và tập dữ liệu}
Nghiên cứu đánh giá các pipeline học máy và Ensemble Learning cho phát hiện xâm nhập mạng trên bộ dữ liệu UNSW-NB15\cite{moustafa2015unsw}. Nhiệm vụ đã được kiểm chứng trong chuỗi thí nghiệm cuối là \textbf{phân loại nhị phân}, với Normal = 0 và Attack = 1; các kết quả theo \texttt{attack\_cat} sau đó chỉ là phân tích nhóm tấn công bị bỏ sót, \emph{không phải mô hình đa lớp}. Mục tiêu là đánh giá performance, lỗi bổ sung, chi phí và độ bất định, không tuyên bố phát minh thuật toán Ensemble Learning mới.

Hai tập official của UNSW-NB15 được giữ nguyên: TRAIN có 175\,341 dòng (56\,000 Normal và 119\,341 Attack), TEST có 82\,332 dòng (37\,000 Normal và 45\,332 Attack). Không gộp rồi chia ngẫu nhiên lại hai tập này. Biến \texttt{id}, \texttt{label} và \texttt{attack\_cat} bị loại khỏi predictors nhị phân; 42 biến còn lại tạo thành track toàn bộ đặc trưng. Bản kiểm toán dữ liệu ghi nhận 8\,541 dòng TEST trùng vector 42 đặc trưng với ít nhất một dòng TRAIN; 73\,791 dòng TEST còn lại tạo \emph{novel-vector sensitivity view}. Đây vẫn là phần của chính official TEST, không phải external, temporal hoặc session split mới.

\section{Feature selection và preprocessing}
Track 17 feature được hình thành từ kết quả lựa chọn trên official TRAIN, cạnh track 42 biến gốc. Năm seed stability \texttt{11, 23, 42, 71, 101} được dùng để kiểm tra tần suất xuất hiện trong từng phương pháp; ngưỡng 0,60 tương ứng ít nhất 3/5 lượt, sau đó giữ biến đạt điều kiện trong ít nhất hai phương pháp. Đây là độ ổn định \emph{của việc chọn biến}, không chứng minh độ ổn định chất lượng dự đoán qua các lần huấn luyện độc lập. Track stable17 đã khóa theo thứ tự:
\begin{quote}\small\raggedright
\texttt{ct\_srv\_dst}, \texttt{ct\_state\_ttl}, \texttt{dload}, \texttt{dttl},
\texttt{proto}, \texttt{state}, \texttt{sttl}, \texttt{ct\_dst\_ltm},
\texttt{dmean}, \texttt{dwin}, \texttt{is\_sm\_ips\_ports},
\texttt{rate}, \texttt{service}, \texttt{swin},
\texttt{ct\_dst\_sport\_ltm}, \texttt{ct\_dst\_src\_ltm}, \texttt{dbytes}.
\end{quote}
Track RFECV riêng tạo 13 biến, không được hoán đổi với stable17. Cả hai kết quả đều nằm trong artifact feature selection đã lưu; không áp dụng lựa chọn đặc trưng mới sau khi xem kết quả TEST. Biến phân loại \texttt{proto}, \texttt{service}, \texttt{state} được One-Hot Encoding với \texttt{handle\_unknown=ignore}; đặc trưng số được đưa qua chuỗi preprocessing theo từng learner, với chuẩn hóa số cho Logistic Regression và không đánh đồng số cột one-hot với 17/42 feature \emph{gốc}. Các phép biến đổi có tham số học được nằm trong \texttt{sklearn.Pipeline}, khớp mã thực nghiệm frozen. Một giới hạn bắt buộc là stable feature selection được thực hiện trên \emph{toàn bộ official TRAIN trước development cross-validation}, không lặp hoàn toàn bên trong từng outer fold. Vì vậy điểm development CV có thể lạc quan hơn kết quả của quy trình feature selection nested thật. Chúng tôi không tuyên bố loại trừ mọi dạng leakage.

\section{Các pipeline được khóa trước final evaluation}
Ba cấu hình cuối và ngưỡng phân loại 0,5 được xác định trước bước đánh giá cuối trên tập TEST.

\begin{table}[htbp]
\centering
\small
\caption{Cấu hình pipeline cuối theo frozen manifest.}
\label{tab:frozen-candidates}
\begin{tabular}{lp{9cm}r}
\hline
Candidate & Cấu hình & Feature \\
\hline
RF-42 & Random Forest 300 cây, \texttt{max\_features}=0.5, \texttt{min\_samples\_leaf}=1 & 42 \\
Stacking-42 & Base Logistic Regression ($C=1$, liblinear), Decision Tree (leaf=2), RF như RF-42; meta Logistic Regression ($C=0.5$, liblinear, class weight None), inner CV 3 folds và \texttt{passthrough}=False & 42 \\
RF-17 & Random Forest 150 cây, \texttt{max\_features}=0.5, \texttt{min\_samples\_leaf}=2 & 17 \\
\hline
\end{tabular}
\end{table}

Stacking sử dụng dự đoán \emph{out-of-fold} của các base learners để huấn luyện meta learner, thay vì sử dụng thuần in-sample predictions của tầng base. Mã nguồn có inner \texttt{StratifiedKFold}; các seed ghi nhận là 42 cho thí nghiệm và 137 cho inner stacking. Việc dùng OOF giúp hạn chế một dạng overfitting tầng meta nhưng phần meta tuning chưa hoàn toàn nested, nên không bảo đảm development validation unbiased. Diversity và voting được xem xét trong development; kết quả voting không thuận lợi cũng là kết quả nghiên cứu cần báo cáo.

RF-17 khác RF-42 cả số feature \emph{và} tham số Random Forest (số cây, kích thước lá), nên mọi chênh lệch giữa chúng là \emph{pipeline-level performance--cost trade-off}, không phải ablation tách biệt nhân quả của số lượng đặc trưng.


\section{Các bước triển khai và huấn luyện}
Quá trình thực nghiệm được thực hiện bằng Python với thư viện \texttt{scikit-learn}. Đối với mỗi cấu hình, chương trình tiến hành thống nhất các bước sau:
\begin{enumerate}
\item Kiểm kê và xác minh hai tập official TRAIN/TEST theo manifest; trong bước chọn feature, development CV và khóa cấu hình, chỉ sử dụng official TRAIN. Cột \texttt{id}, \texttt{attack\_cat}, \texttt{label} không làm predictor nhị phân; file TEST chỉ được đưa vào đánh giá sau freeze.
\item Xây dựng chuỗi tiền xử lý thích hợp với đặc trưng số và đặc trưng phân loại, bảo đảm tham số biến đổi chỉ được ước lượng từ phần dữ liệu huấn luyện được phép.
\item Huấn luyện các mô hình cơ sở; với Stacking, sử dụng dự đoán ngoài fold để xây dựng tập đầu vào cho bộ phân loại tầng kết hợp.
\item Lựa chọn và cố định cấu hình theo kết quả phát triển mô hình trên TRAIN. Sau đó huấn luyện lại bằng toàn bộ TRAIN và dự đoán trên official TEST với ngưỡng 0,5.
\item Tính confusion matrix, Precision, Recall, Macro-F1, MCC, FPR, FNR và thời gian thực thi; xuất dự đoán cùng các bảng chỉ số để phân tích đối chiếu.
\end{enumerate}
Các bước này mô tả đúng trình tự triển khai. Mã nguồn Python đi kèm báo cáo là căn cứ để kiểm tra chi tiết thao tác tiền xử lý, đối tượng mô hình và lệnh đánh giá.

\section{Final evaluation, uncertainty và sai số}
Sau khi đóng băng cấu hình và threshold 0,5, các pipeline được huấn luyện trên official TRAIN và đánh giá cuối trên official TEST. Kết quả dự đoán và các chỉ số đánh giá được trình bày trong Chương~\ref{chap:ket-qua-thao-luan}. Không dùng TEST để chọn lại mô hình, tối ưu ngưỡng hay sửa feature set. Attack là positive class: $\mathrm{FPR}=\mathrm{FP}/(\mathrm{TN}+\mathrm{FP})$, $\mathrm{FNR}=\mathrm{FN}/(\mathrm{TP}+\mathrm{FN})$, Detection Rate = $1-\mathrm{FNR}$. Cùng với Macro-F1 và MCC, nghiên cứu còn lưu ROC-AUC và trường \texttt{pr\_auc}; trường sau được mã nguồn tính bằng \texttt{average\_precision\_score}, vì vậy phải diễn giải là Average Precision (AP), không tự coi là phương pháp tích phân PR-AUC khác.

Chênh lệch metric giữa các model được phân tích bằng paired bootstrap 10\,000 lượt với seed \texttt{20260928}: các dự đoán của hai mô hình trên \emph{cùng dòng} được ghép cặp, phân tầng theo nhãn thật Normal/Attack và giữ nguyên tổng số mẫu của từng lớp bằng resampling multinomial trên các mẫu hình dự đoán liên hợp. CI 95\% là percentile interval cho từng chỉ số riêng, không phải simultaneous familywise-adjusted CI. Đơn vị biến thiên là các dòng benchmark có điều kiện trên số lượng hai lớp, \emph{không} đại diện repeated-training variability, session/host độc lập hay độ bất định ngoài UNSW-NB15. McNemar exact so sánh \emph{paired overall correctness}; hiệu chỉnh Holm thực hiện cho ba cặp model trong từng evaluation view (không phải mọi bảng CI/profile), và đây không phải kiểm định trực tiếp Macro-F1. Các phân tích lỗi theo \texttt{attack\_cat}, \texttt{proto}, \texttt{state}, \texttt{sttl}, \texttt{dload} là mô tả hậu nghiệm; không dùng để suy ra quan hệ nhân quả hay tạo threshold mới.

\section{Các giới hạn về tái lập và tính hợp lệ}
Việc kiểm tra tính toàn vẹn dữ liệu, sự nhất quán giữa cấu hình và mã nguồn, cùng khả năng tính lại chỉ số từ dự đoán giúp tăng tính kiểm chứng của thí nghiệm. Môi trường thư viện của các mô hình lịch sử chưa được lưu đầy đủ, vì vậy chưa chứng minh \emph{đồng nhất bitwise} với đối tượng đã tạo ra số liệu TEST. Tại checkpoint 17/10/2026, cả ba frozen candidates RF-17, RF-42, Stacking-42 đã được \emph{tái dựng riêng} trên 175\,341 dòng official TRAIN; các run reports ghi PASS và kiểm định artifact trên laptop tác giả xác nhận SHA-256, tải lại model, schema và khả năng dự đoán trên 16 dòng TRAIN (69/69 checks). Đây là artifact \texttt{RECONSTRUCTED\_FULL\_TRAIN}, \emph{không phải} model lịch sử và \emph{không} tái chứng minh điểm TEST lịch sử. Kết quả tái dựng có môi trường Python 3.11.9, scikit-learn 1.9.0 được ghi lại. Ranh giới này không được che bởi các thí nghiệm smoke trên tập con.
